\begin{frame}{왜 MLE로 바로 풀리지 않는가: 곱의 합}
  \begin{columns}[T]
    \begin{column}{0.48\textwidth}
      \kb{Week 3 GDA / 나이브베이즈}
      \begin{itemize}\small
        \item ``\textbf{라벨이 주어지면}'' 파라미터별
          likelihood가 \textbf{곱의 합}으로 깔끔하게 풀림
        \item $\Rightarrow$ \textbf{닫힌 형태}의 MLE 공식 존재
      \end{itemize}
    \end{column}
    \begin{column}{0.48\textwidth}
      \kb{GMM}
      \begin{itemize}\small
        \item 라벨 $z$가 \textbf{관측되지 않아}
        \item 전체 likelihood =
          $\prod_i \big(\sum_k \pi_k \mathcal{N}(x_i;
          \mu_k, \sigma_k^2)\big)$
        \item = \textbf{곱의 합}
      \end{itemize}
    \end{column}
  \end{columns}
  \vskip0.8em
  \kq{합이 곱 안에 끼어 있으면 미분을 0으로 놓고 풀 수 있는
  구조가 아니다.}
  \vskip0.6em
  \begin{block}{EM의 본질}
    ``만약 $z$가 \textbf{알렸다면}''은 Week 3 GDA와 \textbf{똑같이}
    닫힌 형태로 풀린다 --- EM은 이 \textbf{``만약''을 반복적으로
    실행}하는 알고리즘.
  \end{block}
\end{frame}
